\section{The non-linear time focused operator}\label{sec_time_focus}

\subsection{Atoms and time focused transform}
\label{sse:definitions.timefocus}
Let us now introduce and study the first non-linear transform of interest here, namely the time focused transform, which involves a focus function defined in the time domain. 


\paragraph{Assumptions}

Throughout this section we will use the following assumptions.

\begin{enumerate}[i.]
    \item
    $h$ is a nonzero, continuous, compactly supported function, called \textit{window}, with length $l\in\R^+_*$. Examples of such a window include most windows used in signal processing (Hann, Blackman,...). 
    \item
    To every $f\in L^2(\R)$ is associated a function $\sigma_f^\tau$, called \textit{focus function}. $\sigma_f^\tau$ will be assumed to be larger than $1$, piecewise continuous and to tend to 1 at infinity, \ie
    \begin{equation}
    \forall f\in L^2(\R), \quad  \sigma^\tau_f \in C^+_{p,0}(\R) + 1.
    \end{equation}
    In addition we will assume that for every $f\in L^2$ there is a sequence $(f_n)_n\in (C_c^\infty)^\N$ such that 
    \begin{equation}\label{eq_hypothesis_cv_sigma}
    f_n \underset{n\rightarrow+\infty}{\overset{L^2}{\longrightarrow}} f \quad \text{ and } \quad \sigma_{f_n}^\tau \underset{n\rightarrow+\infty}{\overset{L^\infty}{\longrightarrow}} \sigma_f^\tau \ .
    \end{equation}
    In order to lighten the notations we will sometimes omit the subscript $\tau$ when there are no ambiguities. 
    \item $\gamma$ is a $C^1$ symmetrical diffeomorphism satisfying $\lim_{t\to -\infty}\gamma(t)=-\infty$  and $\lim_{t\to +\infty}\gamma(t)=+\infty$. 
    %$\ie$ very precise around $0$ (low frequencies) and less precise for high values (high frequencies).
\end{enumerate}


\begin{remark}
    The technical assumption~\eqref{eq_hypothesis_cv_sigma} on the map $f\to\sigma_f^\tau$ is used in the proof of Lemma~\ref{lemma_density_result} below. We believe it should be satisfied for most reasonable choices of mappings $f\to\sigma_f^\tau$.
\end{remark}

\paragraph{Atoms and transform definition} Given the above hypotheses, we can define the  time-focused atoms as
\begin{equation}
\label{fo:time.focused.atoms}
	\forall x,t,\omega \in\R , \ h_{t,\omega,\sigma_f^\tau}(x):=  \sqrt{\gamma'(\omega)\sigma_f^\tau(t)} e^{2i\pi\gamma(\omega)x}h(\sigma_f^\tau(t)(x-t))\ ,
\end{equation}
and the corresponding transform of a given signal $f\in L^2$ by
\begin{equation}\label{eq_def_operator_time_focus}
	\forall t,\omega\in\R, \ M^\tau f(t,\omega) := \langle f , h_{t,\omega,\sigma_f^\tau} \rangle_{L^2}.
\end{equation}
The pointwise definition of the scalar product is guaranteed by the fact that $h(\sigma_f^\tau(t)(x-t))$ is a continuous, compactly supported function of $x$.
When there are no ambiguities, we will sometimes write $h_{t,\omega,f}$ instead of $h_{t,\omega,\sigma_f}$ and $h_{t,\omega,n}$ instead of $h_{t,\omega,\sigma_{f_n}}$ for a certain sequence $(f_n)_n$.
\begin{remark}
\begin{enumerate}
    \item
    In the definition of time-focused atoms~\eqref{fo:time.focused.atoms}, $\sigma_f^\tau(t)$ performs a scaling of the window $h$ around $t$. We stress that the lower bound 1 of $\sigma_f^\tau$ is purely conventional, and states that rescaled windows cannot be larger than $h$. The choice of $h$ and the range of values of $\sigma_f^\tau$ therefore determine the overall resolution of the analysis.
    \item The function $\gamma$ performs a mere relabeling of the frequency axis.
    Examples of $\gamma$ functions include the obvious choice $\gamma(t)=t$, $\gamma$ can also be given a kind of hyperbolic sine shape like, which has the effect of compressing high frequencies.
    %A non-linear $\gamma$ may be used to model a non-linearity in the precision of the analysis along the frequency axis.}
\end{enumerate}
\end{remark}


\paragraph{Density result}

We first prove the following density result, which states that if we take a $L^2$ function $f$ and a $C_c$ function $f_\eps$ as close as possible to $f$ in $L^2$, then the norm of $M^\tau f$ will be controlled by the norm of $M^\tau f_\eps$, which is finite by basic integration rules.

\begin{lemma}\label{lemma_density_result}
    Let $f\in L^2(\R)$ and $(f_n)_{n\in\N}$ be a sequence such that $f_n \in C_c$ for all $n\in\N$, $f_n \rightarrow f$ in $L^2(\R)$ and $\sigma_{f_n}^\tau=:\sigma_n\rightarrow \sigma_f^\tau$ in $L^\infty$. Then
    \begin{equation}
        \forall \eps>0,\ \exists N\in\N,\ \forall n\geqslant N,\ \left|\left\|M^\tau f\right\|_{L^2} - \left\|M^\tau f_n \right\|_{L^2} \right| \leqslant \eps\ .
    \end{equation}
\end{lemma}

\begin{proof}
    Let $\tilde{\eps}>0$ and $n\in\N$ such that $\|f_n-f\|_{L^2}<\tilde{\eps}$ and $\|\sigma_n-\sigma_f\|_\infty<\tilde{\eps}$. We have
    \begin{align*}
        \left| \|M^\tau f_n \|^2_{L^2} - \|M^\tau f\|^2_{L^2} \right| & = \left| \int_{\R^2} |\langle f_n , h_{t,\omega,n}\rangle|^2 d\omega dt - \int_{\R^2} |\langle f , h_{t,\omega,f}\rangle|^2 d\omega dt \right|
    \end{align*}
    We first compute, setting $F_n(x,t):=\sqrt{\sigma_n(t)}f_n(x)h(\sigma_n(t)(x-t))$
    \begin{align*}
       \int_{\R^2} |\langle f_n , h_{t,\omega,n}\rangle|^2 d\omega dt &= 
       \int_{\R^4}F_n(x,t)\bar{F}_n(x',t)\gamma'(\omega)e^{2i\pi\gamma(\omega)(x-x')}dxdx'd\omega dt.
    \end{align*}
    We recognize the formula of $\F^{-1}_1(F_n)(\gamma(u),t)$, hence
    \begin{align*}
        \int_{\R^2} |\langle f_n , h_{t,\omega,n}\rangle|^2 d\omega dt 
        &= \int_{\R^2}\gamma'(\omega)|\F^{-1}_1(F_n)(\gamma(u),t)|^2 d\omega dt \\
        &= \int_{\R^2} |\F^{-1}_1(F_n)(u,t)|^2 dudt \\
       &= \int_{\R^2} |F_n(x,t)|^2 dxdt = \int_{\R^2}\sigma_n(t)|f_n(x)|^2 h^2(\sigma_n(t)(x-t)) dxdt.
    \end{align*}
    The same computations give 
    \begin{equation*}
        \int_{\R^2} |\langle f , h_{t,\omega,f}\rangle|^2 d\omega dt = \int_{\R^2} \sigma_f(t)|f(x)|^2 h^2(\sigma_f(t)(x-t)) dx dt\ .
    \end{equation*}
    We can then write
    \begin{align*}
           \left| \|M^\tau f_n \|^2_{L^2} - \|M^\tau f\|^2_{L^2} \right|  & = \left| \int_\R A(t)dt + \int_\R B(t)dt \right|\ ,
    \end{align*}
    with
    \begin{align*}
        A(t) &:= \int_\R \sigma_f(t)h^2(\sigma_f(t)(x-t))\left(|f_n(x)|^2-|f(x)|^2\right) dx\ .
    \end{align*}
    Therefore,
    \begin{equation*}
        \int_\R |A(t)|dt \leqslant \|\sigma_f\|_\infty \|h\|_{\infty}\tilde{\eps}\ .
    \end{equation*}
    Similarly,
    \begin{equation*}
    B(t):= \int_\R \sigma_n(t)|f_n(x)|^2\left(h^2(\sigma_n(t)(x-t))-h^2(\sigma_f(t)(x-t))\right)dx\ ,
    \end{equation*}
    thus
    \begin{equation*}
        \int_\R |B(t)|dt \leqslant (\|\sigma_f\|_\infty+\tilde{\eps})(\|f\|_{L^2}+\tilde{\eps})^2 2l\,\tilde{\eps}\ .
    \end{equation*}
    Since $\tilde{\eps}$ is arbitrary, the result follows.
\end{proof}

When combined with Proposition~\ref{prop_norm_control_time_operator_L_1_lemma} below, Lemma~\ref{lemma_density_result} will show that $M^\tau$ is well-defined on $L^2$.

\paragraph{Motivations and examples for the time focus function}\label{ssec_time_focus_function}
As stressed in the introduction, it is not the goal of the current paper to discuss in details explicit choices for the focus functions that would be relevant in specific applications. We only provide a couple of prototypical examples, to illustrate desirable and undesirable properties.

Denote by $Vf$ the time-focused transform of $f\in L^2(\R)$, with a constant focus function $\sigma(t)=\sigma_\mathsf{ref}$ (in other words, a STFT with prescribed time-frequency resolution), and assume for simplicity $\gamma(t)=t$. If the goal is to increase the time resolution of the analysis when the analyzed signal $f\in L^2(\R)$ has faster variations, a natural idea could be to consider fixed-time slices of $Lf$ and compute weighted norms of the form
\[
\sigma_f^\tau(t) = 1 + \int_\R w(\omega) |Vf(\omega,t)|\,d\omega \ ,
\]
for some weight function $w$ which enhances the contribution of high frequencies (typically $w(\omega) = |\omega|^n$). This quantity is well-defined as soon as $\sup_t \int_\R w(\omega) |Vf(\omega,t)|\,d\omega<\infty$; furthermore, it defines a continuous function that tends to 1 as $t\to\pm\infty$, since $t\to Vf(t,\omega)$ may be written as the convolution product of two $L^2$ functions (up to a phase factor).

However, such a choice is too naive, as the second term is homogenous of degree 1 with respect to $f$, which results in an increase of the focus function $\sigma$ when $f$ is multiplied by a constant (in which case one can hardly pretend that the resulting function has facter variations). For these reasons, we will privilege non-linear terms, such as terms involving norm ratios. Natural examples involving such ratio are given by entropies such as the Rényi or Shannon entropies
\[
R_\alpha(t) = \frac{1}{1-\alpha} \log \frac{\|Vf(t,\cdot)\|_{L^{\alpha}}^{\alpha}}{\|Vf(t,\cdot)\|_{L^{1}}^{\alpha}}\ ,\qquad H(t) = \lim_{\alpha\to 1} R_\alpha(t)\ ,
\]
assuming the above quantities are well defined and nonzero.
%In this respect, we may notice that $\|Vf(t,\cdot)\|_{L^{2}(d\omega)}^{2} = \int_\R |f(x)|^2|h(x-t)|^2dx$ can only vanish if the window $h$ is compactly supported and $f$ vanishes inside an interval whose length exceeds the support size of $h$ \bruno{...Je sais pas si c'est intéressant de parler de ça...}. 

Entropy is generally used as a measure of spread (or information content as in~\cite{Baraniuk2001measuring}): the more spread out the function, the larger the entropy, independently of global normalization. In the context under consideration here, a large value of $R_\alpha(t)$ would indicate that the "energy" of the spectrogram $|Lf(t,\cdot)|$ is spread throughout the whole frequency domain, which can be interpreted in terms of the presence of a transient event in the signal $f$ at time $t$, and would then require a more time-focused analysis. A corresponding focus function could be defined as
\[
\sigma_f^\tau(t) = A R_\alpha(t) + B\ ,
\]
for some constants $A,B$ that would control the range of $\sigma_f^\tau$, or a similar expression using the Shannon entropy. The mathematical analysis of the behavior of such functions defined on the real line is out of the scope of the present paper. We shall discuss adaptations in discrete, finite-dimensional situations, in Section~\ref{se:numerical.experiments} devoted to numerical illustrations. 
\begin{comment}
\subsection{Example and ideas of the time focus function}\label{ssec_time_focus_function}
Again, it is not the goal of the current paper to discuss in details explicit choices for the focus functions that would be relevant in specific applications. We only provide a simple and generic example that could be of interest (a more sophisticated alternative, developed in the context of audio compression, can be found in~\cite{Molla2004determining}).
\begin{example}[Adaptation to transient events]
\label{ex:time_focus}
As already mentioned in the introduction, a fairly natural idea is to increase the time resolution of the analysis when the analyzed signal has faster variations. This requires a heuristic to measure such speed of variations. As in Example~\ref{ex:frequency_focus}, a ratio of suitably chosen norms can be used for that purpose. Using frequency content as a measure of speed of variations, one may consider the quantity 
\begin{equation}\label{eq_def2_sigma_tau}
	\sigma_{f}^\tau(t) := 1 + \sqrt{\frac{\int_\R |\omega|^{2n} |V_g f(t,\omega)|^2 d\omega}{\int_\R |V_g f(t,\omega)|^2d\omega}}\varphi(t)\ ,
\end{equation}
with $V_gf$ the STFT of $f$ with Gaussian window $g(x):=e^{-x^2/2a^2}$ with variance $a^2$, and $\varphi$ an additional function, which can ensure \pierre{that $\sigma_f^\tau -1$ vanishes at infinity} and appropriate normalization. The rationale is to use the ratio to measure some frequency content of the signal around time $t$ (as an average frequency variable, weighted by the spectrogram at time $t$). The higher the average frequency, the larger the time focus, and smaller the adapted window. A toy example is given in Appendix~\ref{sec:app.time.foc.example} that shows the ability of such functions to detect transient events in a signal and therefore adapt the scale of a window in the time focused transform.

\end{example}
\end{comment}
\subsection{Norm relationship}

We first express the $L^2$ norm of $Mf$ in terms of a certain kernel, which guarantees the well-definedness of the transform on $L^2$, and will be  useful for the rest of the study.

\begin{prop}
    Using the previous definition, for $f\in C_c$ we have then
	\begin{equation}
		\|M^\tau f\|_{L^2(\R^2)}^2 = \int_{\R^2} \hat{f}(\xi)\bar{\hat{f}}(\xi') K_{\sigma_f^\tau}(\xi-\xi') d\xi d\xi',
	\end{equation}
	where the  kernel is 
	\begin{equation}\label{eq_def_time_kernel}
		K_{\sigma_f^\tau}(u) := \int_{\R^2}e^{-2i\pi u t}\bar{\hat{h}}(z)\hat{h}\left(z-\tfrac{u}{\sigma_f^\tau(t)}\right)\, dzdt\ .
	\end{equation}
\end{prop}
\begin{proof}
	In order to make notations lighter we use $\sigma$ for $\sigma_f^\tau$.
	\begin{align*}
		\|M^\tau f\|_{L^2}^2 & = \int_{\R^2} |\langle \hat{f},\hat{h}_{t,\omega,\sigma}\rangle|^2dtd\omega \\
		& = \int_{\R^4} \hat{f}(\xi)\bar{\hat{f}}(\xi')\bar{\hat{h}}_{t,\omega,\sigma}(\xi)\hat{h}_{t,\omega,\sigma}(\xi') d\xi d\xi' dt d\omega \\
		& = \int_{\R^2}\hat{f}(\xi)\bar{\hat{f}}(\xi') \underset{=:K_{\sigma}(\xi,\xi')}{\underbrace{\int_{\R^2}\bar{\hat{h}}_{t,\omega,\sigma}(\xi)\hat{h}_{t,\omega,\sigma}(\xi')dt d\omega}} d\xi d\xi'
	\end{align*}
	We have
	\begin{equation*}
		K_{\sigma}(\xi,\xi'):= \int_{\R^2}\frac{\gamma'(\omega)}{\sigma(t)}e^{-2i\pi(\xi-\xi')t}\bar{\hat{h}}\left(\tfrac{\xi-\gamma(\omega)}{\sigma(t)}\right)\hat{h}\left(\tfrac{\xi'-\gamma(\omega)}{\sigma(t)}\right)\,d\omega dt\ .
	\end{equation*}
	Hence, by setting $y(\omega)=-\frac{\gamma(\omega)}{\sigma(t)}$ we obtain
	\begin{align*}
		K_{\sigma}(\xi,\xi') &= \int_{\R^2} \bar{\hat{h}}\left(\tfrac{\xi}{\sigma(t)}+y\right)\hat{h}\left(\tfrac{\xi'}{\sigma(t)}+y\right)\,e^{-2i\pi(\xi-\xi')t}\,dydt \\
		&= \int_{\R^2}e^{-2i\pi(\xi-\xi')t} \bar{\hat{h}}(z)\hat{h}\left(z-\tfrac{\xi-\xi'}{\sigma(t)}\right)\,dzdt\ .
	\end{align*}
	The last equality holds by applying the changing of variable $z=y+\frac{\xi}{\sigma(t)}$. And now since the kernel is a function of $\xi-\xi'$ we can deduce the expected result.
\end{proof}

\begin{corollary}
	We have, for $f\in C_c(\R)$,
	\begin{equation}
		\|M^\tau f\|_{L^2(\R^2)}^2 = \int_\R |f|^2(t)\F^{-1}(K_{\sigma_f^\tau})(t) dt.
	\end{equation}
\end{corollary}
\begin{proof}
	The proof uses the Plancherel equality and the fact that the Fourier transform of a convolution product is the product of the Fourier transform. It is an elementary computation
	\begin{align*}
		\|M^\tau f\|_{L^2}^2  &= \int_{\R^2} \hat{f}(\xi)\bar{\hat{f}}(\xi') K_{\sigma}(\xi-\xi') d\xi d\xi \\
		& = \int_\R \hat{f}(\xi) \int_\R \bar{\hat{f}}(\xi')K_\sigma(\xi-\xi')d\xi' d\xi \\
		& = \langle \hat{f} , K_\sigma \ast \bar{\hat{f}} \rangle_{L^2_\xi} \\
		& = \int_\R f(t) \F^{-1}\left( K_\sigma \ast \bar{\hat{f}} \right)(t) dt \\
		& = \int_\R |f|^2(t)\F^{-1}(K_\sigma)(t)dt. 
	\end{align*}
\end{proof}

Now we can see that the problem can be solved by controlling the kernel $K_{\sigma_f^\tau}$. In order to obtain such a control, we can introduce the explicit formula of $\F^{-1}(K_{\sigma_f^\tau})$.
\begin{equation}\label{eq_formula_inverse_fourier_kernel}
		\forall t\in\R, \ \F^{-1}(K_{\sigma_f^\tau})(t) = \int_\R \sigma_f^\tau(x)|h(\sigma_f^\tau(x)(x-t))|^2dx.
\end{equation}


\subsection{Main result : norm control}

The rest of the section is dedicated to the proof of Theorem~\ref{theo_bounding_time_marseillan_transform} below.
%This result is the analogous control that we have in section~\ref{sec_frequency_focus} with  Theorem~\ref{theorem_norm_control_frequency_focus}.


\begin{theorem}\label{theo_bounding_time_marseillan_transform}
	Let $f\in L^2(\R)$ and $\sigma_f^\tau$ be a time focus parameter, then
	\begin{equation}
		c_f \|f\|_{L^2(\R)}^2 \leqslant\| M^\tau f\|_{L^2(\R^2)}^2 \leqslant C_f \|f\|_{L^2(\R)}^2\ ,
	\end{equation} 
	where
	\begin{equation}
		c_f = \inf_{t\in\R}\left\{ \int_\R|h(x\sigma_f^\tau(x+t))|^2 dx \right\}>0 \ ,
	\end{equation}
	and
	\begin{equation}
		C_f = \int_\R |h(\sigma_f^\tau(t)t)|^2\sigma_f^\tau(t)dt < \infty \ .
	\end{equation}
\end{theorem}

\begin{remark}
   This result is very similar to a continuous frame condition, except that $c_f$ and $C_f$ depend upon the analyzed function $f$ through the focus function $\sigma_f^\tau$. It is interesting to notice that the non-linearity of the transform only shows up in these constants.
\end{remark}

\subsection{Upper bound control}

In order to prove the upper bound control - which also guarantees the well-definedness in $L^2(\R)\rightarrow L^2(\R^2)$ of our operator - we will control the $L^1$ norm of the previously introduced kernel $K_{\sigma_f^\tau}$. For now, $f$ denotes a $C_c$ function.
\begin{lemma}[$L^1$ norm control]\label{lemma_norm_control_time_kernel_L1}
	Let $\sigma(t)$ be a focus function and $h$ be a window function, both as defined in Section~\ref{sse:definitions.timefocus}. Let the kernel $K_\sigma(u)$ be defined as in equation~\eqref{eq_def_time_kernel}.
%	\begin{equation}
%		K_\sigma(u) = \int_{\R^2} e^{-2i\pi ut}\bar{\hat{h}}(z)\hat{h}(z-\tfrac{u}%{\sigma(t)})dzdt.
%	\end{equation}
	Then
	\begin{equation}
		\int_\R K_\sigma(u)du = \int_{\R} |h(\sigma(t)t)|^2 \sigma(t)dt\ .
	\end{equation}
\end{lemma}
\begin{proof}
	We have by Fubini's Theorem and setting $u' = z - \tfrac{u}{\sigma(t)}$,
	\begin{align*}
		\int_\R K_\sigma(u)du & = \int_{\R^2} \bar{\hat{h}}(z)\int_{\R} e^{-2i\pi ut}\hat{h}(z-\tfrac{u}{\sigma(t)}) du dzdt \\
		& = \int_{\R^2} \bar{\hat{h}}(z) \int_{\R} e^{-2i\pi \sigma(t)t(z-u)}\hat{h}(u) du\sigma(t)dtdz \\
		& = \int_{\R^2} \bar{\hat{h}}(z)e^{-2i\pi \sigma(t)tz}\int_{\R}\hat{h}(u)e^{2i\pi \sigma(t)tu}dudz\sigma(t)dt\\
		& = \int_{\R^2} \bar{\hat{h}}(z)e^{-2i\pi \sigma(t)tz}h(\sigma(t)t) dz\sigma(t)dt \\
		& = \int_{\R} |h(\sigma(t)t)|^2\sigma(t)dt
	\end{align*}
\end{proof}

From Lemma~\ref{lemma_norm_control_time_kernel_L1} we obtain the upper norm control by introducing a weighed window $H_f$ which lightens a bit the notations.
\begin{prop}[Upper bound]\label{prop_norm_control_time_operator_L_1_lemma}
	Let $f\in L^2(\R)$ and introduce the weighed and rescaled window
	\begin{equation}
		H_f(t) := |h(\sigma_f^\tau(t)t)|\sqrt{\sigma_f^\tau(t)}\ ,
	\end{equation}
	then $H_f\in L^2(\R)$ and
	\begin{equation}
		\|Mf\|_{L^2(\R^2)} \leqslant \|f\|_{L^2(\R)}\|H_f\|_{L^2(\R)}.
	\end{equation}
\end{prop}
\begin{proof}
Since $h$ is compactly supported and $\sigma$ bounded and continuous by part, $H_f\in L^2$. Let us now assume $f\in C_c$. For any $f,g\in L^2$ the dual equality
	\begin{equation*}
		\int_{\R} f(t)\hat{g}(t)dt = \int_{\R} \hat{f}(\omega)g(\omega)d\omega
	\end{equation*}
	gives
	\begin{equation*}
		\|Mf\|_{L^2}^2 = \int_{\R} \F(|f|^2)(\omega)K_\sigma(\omega)d\omega\ .
	\end{equation*}
	As $f\in L^2$ we have $\F(|f|^2)\in L^\infty$, hence by using Lemma~\ref{lemma_norm_control_time_kernel_L1} for the second member
	\begin{align*}
		\int_{\R} \F(|f|^2)(\omega)K_\sigma(\omega)d\omega & \leqslant \| \F(|f|^2)\|_{L^\infty}\int_{\R} |h(\sigma(t)t)|^2\sigma(t)dt, \\
		& \leqslant \| |f|^2 \|_{L^1} \| H_f\|_{L^2}^2 \\
		& = \|f\|_{L^2}^2 \|H_f \|_{L^2}^2\ .				
	\end{align*}
	Using the definition of $H_f$ and taking its $L^2$ norm we obtain the bound in Proposition~\ref{prop_norm_control_time_operator_L_1_lemma}. Therefore, the upper bound $C_f$ in Theorem~\ref{theo_bounding_time_marseillan_transform} equals $\|H_f\|_{L^2}$.

 We will now extend the result to $L^2$. Let $f\in L^2$, by assumption~\eqref{eq_hypothesis_cv_sigma} the exists a sequence $(f_n)_n$ with $f_n\in C_c$, converging to $f$ in $L^2$ and such that the sequence $(\sigma_{f_n})_n$ converges to $\sigma_f$ in $L^\infty$. We know by Lemma~\ref{lemma_density_result} that 
 \begin{equation}
     \left|\|M^\tau f\|_{L^2} - \|M^\tau f_n \|_{L^2} \right| \rightarrow 0
 \end{equation}
 However we know by the upper bound control that
 \begin{equation}
     \|M^\tau f_n\|_{L^2} \leqslant \|H_{f_n}\|_{L^2}\|f_n\|_{L^2} \leqslant \sup_n \{\|H_{f_n}\|_{L^2}\}\|f_n\|_{L^2}
 \end{equation}
 Since $\sigma_{f_n}$ converges in $L^\infty$ - by assumption~\eqref{eq_hypothesis_cv_sigma} - we know that $\sup_n \{\|H_{f_n}\|_{L^2}\}< \infty$. Hence, $(\|M^\tau f_n\|)_n$ is bounded, and Lemma~\ref{lemma_density_result} guarantees that the $L^2$ norm of $M^\tau f$ for $f\in L^2$ is finite and controlled by $\|f\|_{L^2}C_f=\|f\|_{L^2}\|H_f\|_{L^2}$.
 
\end{proof}

The following norm control is not necessary at this point but may be useful in future work, so it is presented here.
\begin{prop}[$L^2$ norm control]
	Let $\sigma(t)$ be a focus function and $h$ be a window function, both as defined in Section~\ref{sse:definitions.timefocus}. Let the kernel $K_\sigma(u)$ be defined as in equation~\eqref{eq_def_time_kernel}. Thus we have
	\begin{equation}
		\|K_{\sigma_f^\tau} \|_{L^2(\R)}^2 = \|h\|_{L^2(\R)}^2 \int_{\R} |h(\sigma_f^\tau(t)t)|^2\sigma_f^\tau(t)^2 dt.
	\end{equation}
\end{prop}
\begin{proof}
	Using Fubini's Theorem we obtain
	\begin{equation*}
		\|K_\sigma\|_{L^2}^2 = \int_{\R^4} e^{-2i\pi(u-u')t}|\hat{h}(z)|^2\hat{h}(z-\tfrac{u}{\sigma(t)})\bar{\hat{h}}(z-\tfrac{u'}{\sigma(t)})dudu'dtdz
	\end{equation*}
	We set then $x:=\tfrac{u}{\sigma(t)}$ and $x'=\tfrac{u'}{\sigma(t)}$ which gives
	\begin{align*}
		\|K_\sigma\|_{L^2}^2  &= \int_{\R^2} \underset{(1)}{\underbrace{\left( \int_{\R} e^{-2i\pi x\sigma(t)t}\hat{h}(z-x)dx \right)}}\underset{(2)}{\underbrace{\left( \int_{\R} e^{2i\pi x'\sigma(t)t}\bar{\hat{h}}(z-x')dx' \right)}}|\hat{h}(z)|^2 \sigma^2(t)dtdz.
	\end{align*}
	We set $x_2=z-x$ and hence we recognize the inverse Fourier transform formula
	\begin{align*}
		(1) &= \int_{\R} e^{-2i\pi(z-x_2)\sigma(t)t}\hat{h}(x_2)dx_2 \\
		& = e^{-2i\pi z\sigma(t)t}h(\sigma(t)t).
	\end{align*}
	Furthermore $(2) = \overline{(1)} $. Thus, when we gather all the members we obtain by Fourier isometry
	\begin{align*}
		\|K_\sigma\|_{L^2}^2 &= \int_{\R^2} |h(\sigma(t)t)|^2 |\hat{h}(z)|^2 dz \sigma(t)^2dt \\
		& = \|h\|_{L^2}^2 \int_{\R} |h(\sigma(t)t)|^2\sigma(t)^2 dt.
	\end{align*}
\end{proof}




\subsection{Lower bound control}

The following proposition gives a strictly positive lower bound for the norm, in the case of non zero signals.
\begin{prop}[Lower bound]\label{prop_lower_bound_control_time_variance}
	Let $\sigma(t)$ be a focus function and $h$ be a compactly supported window function, both as defined in Section~\ref{sse:definitions.timefocus}. There exists $c=c(\sigma)>0$ which depends on $\sigma$ such that
	\begin{equation}\label{eq_lower_bound_kernel}
		\forall t\in\R, \ \F^{-1}(K_\sigma)(t) > c\ ,
	\end{equation}
	with
    \begin{equation}
		c:= \inf_{t\in\R}\left\{ \int_\R|h(x\sigma(x+t))|^2 dx \right\} .
	\end{equation}
\end{prop}
\begin{proof}
Let $t\in\R$. Since $\sigma(x)\geqslant1$ for any $x\in\R$, by using equation~\eqref{eq_formula_inverse_fourier_kernel} we have
\begin{align*}
    \mathcal{F}^{-1}(K_\sigma)(t) & \geqslant \int_\R |h(\sigma(x)(x-t))|^2 dx
\end{align*}
Hence we have
\begin{equation}
    \forall t\in\R, \ \mathcal{F}^{-1}(K_\sigma)(t)\geqslant H(t).
\end{equation}
Now let us prove that $\inf_t H(t) =: c >0$. We obviously have $H(t)>0$ for any $t\in\R$ furthermore by the use of the dominated convergence theorem we also have $\lim_{\pm\infty}H(t)>0$. And since $H(t)$ is continuous (again by the dominated convergence theorem) we can conclude that there exists $c=c(\sigma)>0$ such that
\begin{equation}
    \inf_{t\in\R}H(t)=c>0\ .
\end{equation}
\end{proof}

\begin{remark}
Since $h$ is supposed continuous (and nonzero), a lower bound independent of $\sigma$ can be obtained for $c$. Without loss of generality, assume that $|h|$ attains its maximum value $\|h\|_\infty$ at the origin. Then there exists $a>0$ such that for every $y\in (-a,a)$, $|h(y)|>  \|h\|_\infty/\sqrt{2}$. Since $\sigma(x)\geqslant 1$ for all $x$, we have $x\in (t-a/\sigma(x),t+a/\sigma(x))$ for every $x\in (t-a,t+a)$ so that $(x-t)\sigma(x)\in (-a,a)$. Therefore, we may write
\[
\int_\R |h(x\sigma(x+t))|^2dx = \int_\R |h((x-t)\sigma(x)|^2dx\ge \int_{(t-a,t+a)}|h((x-t)\sigma(x)|^2dx> a \|h\|_\infty^2\ ,
\]
which doesn't depend on $\sigma$, then on $f$ if $\sigma=\sigma_f^\tau$.

Note that the compact support assumption is not necessary for that lower bound.
\end{remark}
